\chapter{Lười, trì hoãn và ảo tưởng cố gắng}
\markboth{Lười, trì hoãn và ảo tưởng cố gắng}{Lười, trì hoãn và ảo tưởng cố gắng}

\section{Tối nay em định học mà cầm điện thoại một chút thôi}
\begin{truyen}{Tối nay em định học mà cầm điện thoại một chút thôi}{Classroom Talk}
\chuhoa{M}{ột bạn thú nhận không vòng vo: ``Em mở điện thoại tính coi 5 phút nghỉ đầu óc thôi. Xong ngẩng lên là mất gần tiếng.''}

Thầy hỏi: ``Vậy em bị thiếu thời gian hay thiếu thắng nổi cái tiện trước mắt?''

Bạn ấy cười trừ: ``Chắc cái thứ hai.''

Thầy đáp: ``Đúng. Nhiều khi mình không lười kiểu nằm im. Mình lười kiểu tự cho phép mình trượt dần.''

\textit{(Procrastination often does not feel like laziness at first. It feels like a small harmless delay that quietly grows.)}
\end{truyen}

\begin{baihoc}
Nhiều buổi học hỏng không vì thiếu quyết tâm lớn, mà vì thua mấy cái tiện nhỏ lặp đi lặp lại.
\textit{(Many study sessions fail not from a lack of grand determination, but from repeated surrender to small comforts.)}
\end{baihoc}

\begin{ghinhoanh}
\textbf{Short English to remember:}

Small distractions become big losses.

Delay grows when you keep permitting it.
\end{ghinhoanh}

\begin{tuhocnhanh}
1. Cái gì hay kéo em trượt khỏi bàn học nhất? \textit{(What most often pulls you away from study?)}

2. Em có cách nào chặn nó trước khi bắt đầu học không? \textit{(Can you block it before your study session starts?)}
\end{tuhocnhanh}
\ngancach

\section{Em làm bài muộn chứ không phải em không làm}
\begin{truyen}{Em làm bài muộn chứ không phải em không làm}{Classroom Talk}
\chuhoa{C}{ó bạn chống chế rất quen: ``Em vẫn làm mà thầy. Chỉ là em hay làm sát giờ thôi, chứ em đâu bỏ.''}

Thầy hỏi: ``Làm sát giờ rồi nộp thiếu, nộp ẩu, thức khuya vật vờ hôm sau, vậy khác bỏ bao nhiêu?''

Bạn ấy cứng họng mấy giây rồi nói: ``Nghe hơi đau.''

Thầy đáp: ``Đau vì đúng. Nhiều người không bỏ việc, chỉ bỏ chất lượng và bỏ sức khỏe trước.''

\textit{(Doing everything at the last minute is not responsibility. It often means sacrificing quality and health while pretending the task was still done.)}
\end{truyen}

\begin{baihoc}
Trì hoãn lâu ngày làm mình tưởng mình vẫn ổn, cho tới khi kết quả và cơ thể cùng phản ứng.
\textit{(Procrastination can make you feel functional until both your results and your body start objecting.)}
\end{baihoc}

\begin{ghinhoanh}
\textbf{Short English to remember:}

Late work is not always real discipline.

Finishing badly is not the same as finishing well.
\end{ghinhoanh}

\begin{tuhocnhanh}
1. Em hay nộp đúng hạn hay hay tự cứu phút cuối? \textit{(Do you submit on time, or survive at the last second?)}

2. Trì hoãn đang làm hỏng điều gì ngoài điểm số? \textit{(What besides grades is procrastination damaging?)}
\end{tuhocnhanh}
\ngancach

\section{Em học theo cảm hứng nên mất nhịp suốt}
\begin{truyen}{Em học theo cảm hứng nên mất nhịp suốt}{Classroom Talk}
\chuhoa{M}{ột bạn nói rất thành thật: ``Lúc có mood em học sung lắm. Nhưng không có hứng là em tắt luôn. Em sống kiểu đó riết nên lúc lên lúc xuống.''}

Thầy cười: ``Nếu đợi hứng mới làm, sau này đi làm em cũng toang với đời.''

Bạn ấy hỏi: ``Sao nặng vậy thầy?''

Thầy nói: ``Vì việc quan trọng hiếm khi chờ cảm hứng. Người bền là người làm được cả lúc không thích.''

\textit{(Motivation rises and falls. Stable progress comes from acting even when excitement is absent.)}
\end{truyen}

\begin{baihoc}
Cảm hứng giúp bắt đầu nhanh. Kỷ luật mới giúp không bỏ dở giữa đường.
\textit{(Motivation helps you start. Discipline keeps you from collapsing halfway.)}
\end{baihoc}

\begin{ghinhoanh}
\textbf{Short English to remember:}

Motivation is helpful.

Discipline is dependable.
\end{ghinhoanh}

\begin{tuhocnhanh}
1. Lịch học của em có phụ thuộc tâm trạng quá nhiều không? \textit{(Does your study routine depend too much on mood?)}

2. Em có một thói quen nhỏ nào làm được dù không có hứng? \textit{(What small habit can you still do without motivation?)}
\end{tuhocnhanh}
\ngancach

\section{Em muốn giỏi nhanh, giỏi chậm nản lắm}
\begin{truyen}{Em muốn giỏi nhanh, giỏi chậm nản lắm}{Classroom Talk}
\chuhoa{C}{ó bạn than: ``Em ghét cảm giác tập hoài mà chưa hay. Em muốn thấy tiến bộ lẹ, chứ chậm quá em tụt mood.''}

Thầy hỏi: ``Em muốn giỏi thật hay muốn cảm giác mình giỏi nhanh?''

Bạn ấy im rồi cười: ``Chắc em mê cái cảm giác hơn.''

Thầy gật đầu: ``Nhiều người bỏ cuộc không phải vì không tiến bộ, mà vì tiến bộ không kịp làm họ sướng.''

\textit{(The hunger for fast improvement can become its own trap. Real mastery often grows slower than ego wants.)}
\end{truyen}

\begin{baihoc}
Người bền là người chịu được giai đoạn chưa giỏi mà vẫn làm tiếp.
\textit{(A durable learner can tolerate being unskilled for a while and still continue.)}
\end{baihoc}

\begin{ghinhoanh}
\textbf{Short English to remember:}

Fast progress feels good.

Slow mastery lasts longer.
\end{ghinhoanh}

\begin{tuhocnhanh}
1. Em đang theo đuổi năng lực thật hay cảm giác thắng nhanh? \textit{(Are you chasing real ability or quick emotional reward?)}

2. Em có chấp nhận một quá trình chậm nhưng chắc không? \textit{(Can you accept a slow but solid process?)}
\end{tuhocnhanh}
\ngancach